\documentclass[../main]{subfiles}
\begin{document}
\chapter{Resistencia Equivalente}
\img{images/01/resistencias.jpg}{Resistencias eléctricas.}
\info{
Resistencia Equivalente}{
  \subsubsection{Paralelo}
  \emph{La resistencia equivalente es la resultante total presente
    resistencias de un circuito, el valor total de las resistencias en
    paralelo es igual al recíproco de la suma de los recíprocos de las
  resistencias.}

  \begin{equation*}
    R_T= \frac{1}{\frac{1}{R_1}+\frac{1}{R_2}} \Rightarrow
    R_T= \frac{1}{\frac{1}{R_1}+\frac{1}{R_2}+\dots+\frac{1}{R_n}}
  \end{equation*}

  \subsubsection{Serie}
  \emph{La resistencia equivalente de resistores conectados en serie
  es igual a la suma de las resistencias individuales.}

  \begin{equation*}
    R_T= R_1+R_2 \Rightarrow R_T= R_1+R_2+\dots+R_n
  \end{equation*}

  Equivalencias: La unidad para resistencias es el Ohm $\Omega$.

}
\actividad{Calcular la resistencia equivalente entre $a$ y $b$.}

$R_1=2 \Omega$
\quad
$R_2=4 \Omega$
\quad
$R_3=6 \Omega$
\quad
$R_4= 5 \Omega$
\quad
$R_5 = 6 \Omega$
\quad
$R_6 = 3 \Omega$
\quad
$R_7= 7 \Omega$

\begin{enumerate}
  \item.
    \begin{center}
      \begin{circuitikz}
        \draw (-2,0) node[label={above:a}]{} to[resistor=$R_1$,,*-] (0,0);
        \draw (0,0) to[resistor={$R_2$}] (2,0);
        \draw (2,0) to[resistor=$R_3$,,-*] (4,0) node[label={above:b}]{};

      \end{circuitikz}
    \end{center}
  \item .
    \begin{center}
      \begin{circuitikz}
        \draw (-3,0) node[label={above:a}]{}
        to[resistor=$R_2$,*-*](3,0) node[label={above:b}]{}
        (2,0)--(2,2) to[resistor=$R_1$] (-2,2)--(-2,0);
        \draw (-2,0)--(-2,-2);
        \draw (2,0)--(2,-2);
        \draw (-2,-2) to[resistor=$R_3$] (2,-2);
      \end{circuitikz}
    \end{center}
  \item .
    \begin{center}
      \begin{circuitikz}
        \draw (-3,0) node[label={above:a}]{}
        to[resistor=$R_1$,*-](0,0) to[resistor=$R_2$,-*] (3,0)
        node[label={above:b}]{} (2,0)--(2,-2) to[resistor=$R_3$]
        (-2,-2)--(-2,0);
      \end{circuitikz}
    \end{center}
  \item .
    \begin{center}
      \begin{circuitikz}
        \draw (-3,0) node[label={above:a}]{}
        to[resistor=$R_1$,*-*](0,3) to[resistor=$R_2$,-*] (3,0)
        node[label={above:b}]{}  to[resistor=$R_3$] (-3,0);
      \end{circuitikz}
    \end{center}
  \item .
    \begin{center}
      \begin{circuitikz}
        \draw (-3,0) node[label={above:a}]{}
        to[resistor=$R_1$,*-*](0,3) to[resistor=$R_4$,-*] (3,0)
        node[label={above:b}]{}  to[resistor=$R_5$] (-3,0);
      \end{circuitikz}
    \end{center}
  \item .
    \begin{center}
      \begin{circuitikz}
        \draw (-3,0)  to[resistor=$R_1$,*-*](0,0)
        node[label={above:a}]{} to[resistor=$R_2$,-*] (3,0)
        node[label={above:b}]{} (2,0)--(2,-2) to[resistor=$R_3$]
        (-2,-2)--(-2,0);

      \end{circuitikz}
    \end{center}
  \item .
    \begin{center}
      \begin{circuitikz}
        \draw (-3,0) node[label={above:a}]{}
        to[resistor=$R_2$,*-*](2,0)  (2,0)--(2,2) to[resistor=$R_1$]
        (-2,2)--(-2,0);
        \draw (-2,0)--(-2,-2);
        \draw (2,0)--(2,-2);
        \draw (-2,-2) to[resistor=$R_3$] (2,-2);
        \draw (2,0) to[resistor=$R_4$] (4,0);
        \draw (4,0) to[resistor=$R_5$,-*] (6,0) node[label={above:b}]{};
      \end{circuitikz}
    \end{center}
  \item .
    \begin{center}
      \begin{circuitikz}
        \draw (-3,0) node[label={above:a}]{}
        to[resistor=$R_1$,*-](3,0) to[resistor=$R_2$,-*] (3,-3)
        node[label={right:b}]{};
        \draw (-3,0)  to[resistor=$R_3$,*-](-3,-3) to[resistor=$R_4$,-*] (3,-3);
      \end{circuitikz}
    \end{center}
  \item .
    \begin{center}
      \begin{circuitikz}
        \draw (-3,0)node[label={above:a}]{} -- (-2,1)
        to[resistor=$R_1$,*-*](0,1)
        to[resistor=$R_2$,*-*] (2,-1)
        to[resistor=$R_3$,*-*] (0,-3)
        to[resistor=$R_4$,*-*] (-2,-3)
        -- (-2.5,-2.5)--(-2,-2)
        to[resistor=$R_5$,*-*] (-3,-1)--(-3.5,-1.5)--(-4,-1)
        node[label={above:b}]{};
        \draw (-2.5,-2.5)--(-3,-3)
        to[resistor=$R_6$,*-*] (-4,-2) --(-3.5,-1.5);
      \end{circuitikz}
    \end{center}
  \item .
    \begin{center}
      \begin{circuitikz}
        \draw (-3,0) node[label={above:a}]{}
        to[resistor=$R_1$,*-](3,0) to[resistor=$R_4$,-*] (3,-3)
        node[label={right:b}]{};
        \draw (-3,0)  to[resistor=$R_2$,*-](-3,-3) to[resistor=$R_5$,-*] (3,-3);
        \draw (-3,0) to[resistor=$R_3$,-*] (3,-3);
      \end{circuitikz}
    \end{center}
  \item .
    \begin{center}
      \begin{circuitikz}
        \draw (0,2) node[label={above:a}]{} to[resistor=$R_1$,*-*](2,4)
        to[resistor=$R_2$,-*] (4,4)
        to[resistor=$R_3$,-*] (6,2)
        to[resistor=$R_4$,-] (0,2)
        to[resistor=$R_5$,-*] (2,0)
        to[resistor=$R_6$,-*] (4,0)
        to[resistor=$R_7$,-] (6,2) node[label={above:b}]{};
      \end{circuitikz}
    \end{center}
  \item .
    \begin{center}
      \begin{circuitikz}
        \draw (0,1)--(1,0)
        to[resistor=$R_5$,*-*] (4,0) node[label={right:b}]{}
        to[resistor=$R_3$,*-*] (4,3)node[label={right:a}]{} --(3,4)
        to[resistor=$R_2$,*-*] (0,4)
        to[resistor=$R_4$,-*] (0,1);
        \draw (0,4) to[resistor=$R_1$,*-*] (4,0);

      \end{circuitikz}
    \end{center}
\end{enumerate}
\end{document}

